% !PRINTED 513
\weberpage{513}
%<<CONT
den Gruppen $\pi^{-1}Q\pi$ aller dieser Functionen, also auch in ihrem
gr\"ossten gemeinsamen Theiler $R$ vorkommen. Die Operationen
von $R$ m\"ogen mit $\sigma$ bezeichnet sein. Sind dann $\pi$ und $\pi'$ zwei
Operationen, die unter den Gr\"ossen (1) dieselbe Permutation
hervorrufen, so wird $\pi'\pi^{-1}$ die urspr\"ungliche Anordnung der $\psi$
wieder herstellen, also gleich einer der Gr\"ossen $\sigma$ sein, oder
\[ \pi' = \sigma\pi. \]
Wir haben daher das Ergebniss:

\begin{quote}
\gesperrt{Die Permutationen der Nebengruppe $R\pi$ und nur
diese rufen unter den Gr\"ossen (1) eine und dieselbe
Permutation hervor.}
\end{quote}

Der Grad der \gesperrt{Galois}'schen Gruppe der Resolvente $\varphi(t) = 0$
ist also gleich der Anzahl dieser Nebengruppen, d. h. gleich dem
Quotienten $p : r$ oder dem Index des Theilers $R$ von $P$.

Ist $R$ die identische Gruppe, also $r = 1$, und folglich $\varphi(t) = 0$
eine Totalresolvente, so ist der Grad ihrer Gruppe ebenso hoch,
wie der Grad der Gruppe der urspr\"unglichen Gleichung, und beide
Gruppen sind \"uberdies isomorph. In Bezug auf die Gruppe ist
also nichts gewonnen. Die gegebene Gleichung ist mit der
Resolvente, so verschieden auch ihre Grade sein m\"ogen, \"aquivalent.

Ist auf der anderen Seite $Q$ ein Normaltheiler, also $R$ mit $Q$
identisch, so ist $\varphi(t) = 0$ eine Partialresolvente und der Grad
ihrer Gruppe ist gleich dem Index $j$ des Theilers $Q$ von $P$.
Nach der Adjunction einer Wurzel dieser Resolvente reducirt
sich die Gruppe der urspr\"unglichen Gleichung auf $Q$, also auf
den Grad $q$. Es ist also eine Spaltung der Gruppe erfolgt.

Wenn man Resolventen bilden will von m\"oglichst niedrigem
Grade, so hat man Theiler der Gruppe $P$ aufzusuchen von m\"oglichst kleinem Index, also von m\"oglichst hohem Grade; soll aber
gleichzeitig eine Reduction der Gruppe eintreten, so m\"ussen die
Theiler Normaltheiler sein.

\begin{center}
\S.\,157.

\gesperrt{Reduction der Galois'schen Gruppe durch Adjunction
beliebiger Irrationalit\"aten.}
\end{center}

Die Aufgabe der L\"osung einer algebraischen Gleichung wird
nach der \gesperrt{Galois}'schen Auffassung durch eine andere ersetzt,
%CONT>>
